% GENERATED FILE. Edit canonical lessons, then run npm run generate:books.
% canonical-source-hash: fnv1a64:5b92f2c5fafd9d8e
% canonical-lessons: UR-C274-agar-to, UR-C274-agar-barish-hui, UR-C274-shart, UR-R274-if-and-otherwise

\chapter{If, Then, or Else}
\label{ch:ur-if-then-or-else}
\hlchaptermodality{274}

\begin{chapteropening}
\textbf{By the end of this chapter:} \emph{I can say what will happen if something happens (agar bārish huī, to …; agar voh āe gī, to …), give a warning with varnā, and set a condition (ek shart par).}

Complete UR-R274-if-and-otherwise: say what will happen if something happens, or else.
\end{chapteropening}

\section[\texorpdfstring{\ur{اگر} … \ur{تو} (agar … to)}{اگر … تو (agar … to)}]{\ur{اگر} … \ur{تو} (agar … to) — if, then}

\label{lesson:UR-C274-agar-to}



\begin{quote}
Say \textquotedblleft{}if\textquotedblright{}. (\emph{agar}.) Then \textquotedblleft{}then I'll stay at home\textquotedblright{}. (\emph{to maiṅ ghar par rahūṅ gā}.) \textbf{\ur{اگر}} opens a condition; \textbf{\ur{تو}} opens what follows.
\end{quote}

\subsection*{You'll want to know: \ur{اگر} … \ur{تو}}
\begin{quote}
\textbf{\ur{اگر} \ur{وقت} \ur{ہے،}}
\end{quote}

\begin{quote}
\emph{agar vaqt hai,}
\end{quote}

\begin{quote}
\textbf{\ur{تو} \ur{آئیے۔}}
\end{quote}

\begin{quote}
\emph{to āiye.} — \textquotedblleft{}If you have time, come.\textquotedblright{}
\end{quote}

\textbf{\ur{اگر}} (\emph{agar}), \textquotedblleft{}if\textquotedblright{}, opens the condition; \textbf{\ur{تو}} (\emph{to}), \textquotedblleft{}then\textquotedblright{}, opens the result. In talk, \textbf{\ur{تو}} often drops out.

\begin{quote}
\textbf{\ur{اگر} \ur{آپ} \ur{چاہتے} \ur{ہیں،}}
\end{quote}

\begin{quote}
\emph{agar āp chāhte haiṅ,}
\end{quote}

\begin{quote}
\textbf{\ur{تو} \ur{ہم} \ur{چلیں} \ur{گے۔}}
\end{quote}

\begin{quote}
\emph{to ham chaleṅ ge.} — \textquotedblleft{}If you like, we'll go.\textquotedblright{}
\end{quote}

\subsection*{You'll want to know: \ur{ورنہ}, or else}
\begin{quote}
\textbf{\ur{جلدی} \ur{چلیے،}}
\end{quote}

\begin{quote}
\emph{jaldī chaliye,}
\end{quote}

\begin{quote}
\textbf{\ur{ورنہ} \ur{گاڑی} \ur{چھوٹ} \ur{جائے} \ur{گی۔}}
\end{quote}

\begin{quote}
\emph{varnā gāṛī chhūṭ jāe gī.} — \textquotedblleft{}Hurry, or we'll miss the train.\textquotedblright{}
\end{quote}

\textbf{\ur{ورنہ}} (\emph{varnā}), \textquotedblleft{}otherwise\textquotedblright{}, turns the condition round: do this, or that will happen.

\subsection*{Your turn}
Say these aloud:

\begin{itemize}
\raggedright
  \item \textquotedblleft{}If you have time, come\textquotedblright{} — \emph{agar vaqt hai, to āiye}
  \item \textquotedblleft{}If you have time, sit down\textquotedblright{} — \emph{agar vaqt hai, to baiṭhiye}
  \item \textquotedblleft{}If you have time, stay\textquotedblright{} — \emph{agar vaqt hai, to rukiye}
\end{itemize}

\begin{tcolorbox}[breakable,colback=teal!4,colframe=teal!35!black,title={Before you move on}]
Which word opens the condition? (\textbf{\ur{اگر}}.) Which opens the result? (\textbf{\ur{تو}}.) \textquotedblleft{}Or else\textquotedblright{}? (\textbf{\ur{ورنہ}}.)
\end{tcolorbox}
\section[\texorpdfstring{\ur{اگر} \ur{بارش} \ur{ہوئی} (agar bārish huī)}{اگر بارش ہوئی (agar bārish huī)}]{\ur{اگر} \ur{بارش} \ur{ہوئی} (agar bārish huī) — if it rains}

\label{lesson:UR-C274-agar-barish-hui}



\begin{quote}
Say \textquotedblleft{}it will rain\textquotedblright{}. (\emph{bārish ho gī}.) Then \textquotedblleft{}if you have time, come\textquotedblright{}. (\emph{agar vaqt hai, to āiye}.) Now a condition about tomorrow.
\end{quote}

\subsection*{You'll want to know: the plain past after \ur{اگر}}
\begin{quote}
\textbf{\ur{اگر} \ur{بارش} \ur{ہوئی،}}
\end{quote}

\begin{quote}
\emph{agar bārish huī,}
\end{quote}

\begin{quote}
\textbf{\ur{تو} \ur{ہم} \ur{گھر} \ur{پر} \ur{رہیں} \ur{گے۔}}
\end{quote}

\begin{quote}
\emph{to ham ghar par raheṅ ge.} — \textquotedblleft{}If it rains, we'll stay at home.\textquotedblright{}
\end{quote}

\textbf{\ur{ہوئی}} (\emph{huī}) is the plain past, but the rain has not come: Urdu pictures it as already there, and the result is in the future.

\begin{quote}
\textbf{\ur{اگر} \ur{گاڑی} \ur{چھوٹ} \ur{گئی،}}
\end{quote}

\begin{quote}
\emph{agar gāṛī chhūṭ gaī,}
\end{quote}

\begin{quote}
\textbf{\ur{تو} \ur{ہم} \ur{بس} \ur{سے} \ur{جائیں} \ur{گے۔}}
\end{quote}

\begin{quote}
\emph{to ham bas se jāeṅ ge.} — \textquotedblleft{}If we miss the train, we'll go by bus.\textquotedblright{}
\end{quote}

\subsection*{You'll want to know: the future after \ur{اگر}}
\begin{quote}
\textbf{\ur{اگر} \ur{وہ} \ur{کل} \ur{آئے} \ur{گی،}}
\end{quote}

\begin{quote}
\emph{agar voh kal āe gī,}
\end{quote}

\begin{quote}
\textbf{\ur{تو} \ur{ہم} \ur{بازار} \ur{جائیں} \ur{گے۔}}
\end{quote}

\begin{quote}
\emph{to ham bāzār jāeṅ ge.} — \textquotedblleft{}If she comes tomorrow, we'll go to the market.\textquotedblright{}
\end{quote}

The future can stand in both halves. Either way, the result is what will happen.

\subsection*{Your turn}
Say these aloud:

\begin{itemize}
\raggedright
  \item \textquotedblleft{}If it rains, we'll stay at home\textquotedblright{} — \emph{agar bārish huī, to ham ghar par raheṅ ge}
  \item \textquotedblleft{}If it rains, we'll go by bus\textquotedblright{} — \emph{agar bārish huī, to ham bas se jāeṅ ge}
  \item \textquotedblleft{}If it rains, we'll take a taxi\textquotedblright{} — \emph{agar bārish huī, to ham ṭaiksī leṅ ge}
\end{itemize}

\begin{tcolorbox}[breakable,colback=teal!4,colframe=teal!35!black,title={Before you move on}]
\textquotedblleft{}If it rains\textquotedblright{}? (\textbf{\ur{اگر} \ur{بارش} \ur{ہوئی}}.) Has it rained? (No: Urdu pictures it.) \textquotedblleft{}If she comes tomorrow\textquotedblright{}? (\textbf{\ur{اگر} \ur{وہ} \ur{کل} \ur{آئے} \ur{گی}}.)
\end{tcolorbox}
\section[\texorpdfstring{\ur{شرط} (shart)}{شرط (shart)}]{\ur{شرط} (shart) — a condition}

\label{lesson:UR-C274-shart}



\begin{quote}
Before the new one: say \textquotedblleft{}if she comes tomorrow, we'll go\textquotedblright{}. (\emph{agar voh kal āe gī, to ham jāeṅ ge}.) Her coming is the condition.
\end{quote}

\subsection*{You'll want to know: \ur{شرط}}
\textbf{\ur{شرط}} — \emph{shart} — \textquotedblleft{}a condition\textquotedblright{}. It is feminine.

\begin{quote}
\textbf{\ur{ایک} \ur{شرط} \ur{ہے۔}}
\end{quote}

\begin{quote}
\emph{ek shart hai.} — \textquotedblleft{}There's one condition.\textquotedblright{}
\end{quote}

\begin{quote}
\textbf{\ur{میں} \ur{ایک} \ur{شرط} \ur{پر} \ur{آؤں} \ur{گا۔}}
\end{quote}

\begin{quote}
\emph{maiṅ ek shart par āūṅ gā.} — \textquotedblleft{}I'll come on one condition.\textquotedblright{}
\end{quote}

\begin{quote}
\textbf{\ur{شرط} \ur{یہ} \ur{ہے} \ur{کہ} \ur{آپ} \ur{بھی} \ur{آئیں} \ur{گے۔}}
\end{quote}

\begin{quote}
\emph{shart yih hai ki āp bhī āeṅ ge.} — \textquotedblleft{}The condition is that you'll come too.\textquotedblright{}
\end{quote}

\begin{cousinweb}
\textbf{\ur{شرط}} (\emph{shart}) is usually traced, through Persian, to Arabic \emph{sharṭ}, \textquotedblleft{}a condition\textquotedblright{}.
\end{cousinweb}

\subsection*{Your turn}
Say these aloud:

\begin{itemize}
\raggedright
  \item \emph{shart}, then \emph{ek shart par}
  \item \textquotedblleft{}I'll go on one condition\textquotedblright{} — \emph{maiṅ ek shart par jāūṅ gā}
\end{itemize}

\begin{tcolorbox}[breakable,colback=teal!4,colframe=teal!35!black,title={Before you move on}]
What does \textbf{\ur{شرط}} mean? (\textquotedblleft{}A condition\textquotedblright{}.) \textquotedblleft{}On one condition\textquotedblright{}? (\textbf{\ur{ایک} \ur{شرط} \ur{پر}}.)
\end{tcolorbox}
\section[(review)]{Review — if, then, or else}

\label{lesson:UR-R274-if-and-otherwise}



\begin{quote}
Say \textquotedblleft{}the bus is cheaper than a taxi\textquotedblright{} and \textquotedblleft{}this is the best\textquotedblright{}. (\emph{bas ṭaiksī se sastī hai}, \emph{yih sab se achchhā hai}.) Then \textquotedblleft{}I've lived here for two years\textquotedblright{}. (\emph{maiṅ do sāl se yahāṅ rahtā hūṅ}.)
\end{quote}

\subsection*{Your turn}
Say these aloud:

\begin{itemize}
\raggedright
  \item I hope she'll come. If she comes tomorrow, we'll go to the market.
  \item If you have time, come.
  \item Hurry, or we'll miss the train.
  \item If it rains, we'll stay at home.
  \item If we miss the train, we'll go by bus.
  \item I'll come on one condition.
\end{itemize}

\begin{tcolorbox}[breakable,colback=teal!4,colframe=teal!35!black,title={Before you move on}]
\textquotedblleft{}If it rains\textquotedblright{}? (\textbf{\ur{اگر} \ur{بارش} \ur{ہوئی}}.) Now the other way round: what does \textbf{\ur{ورنہ}} warn about? (What will happen if you don't.) In \textbf{\ur{اگر} \ur{وہ} \ur{کل} \ur{آئے} \ur{گی}}, which time is the verb in? (The future.)
\end{tcolorbox}
